% ===================================================================
%  نقشہ نمبر 1 — سکول۔ — ثانوی سکول۔ — جماعت۔
%
%  Ceci n'est PAS une transcription : c'est une traduction, faite en
%  2026, en pendjabi occidental d'aujourd'hui, ecrit en SHAHMOUKHI.
%  D'ou trois differences avec texto/io et texto/fr, et trois
%  seulement :
%
%   * pas de \nl ni de \cc — la traduction ne suit aucune ligne
%     imprimee, et les lignes de ce fichier ne sont que du confort ;
%   * pas de \begin{VUpage} — elle n'a ni feuillet ni folio, et la
%     page de lecture ne lui en affiche pas ;
%   * le gras se pose sur le TERME seul, ponctuation dehors.
%
%  Tout le reste est identique, et doit l'etre : les cles « %%K » sont
%  celles de texto/io, macro pour macro pour l'apparat, et les renvois
%  \textsuperscript{(n)} visent les MEMES objets de la planche — ce
%  sont eux qui ouvrent les gros plans.
%
%  DEUX COLONNES POUR UNE SEULE LANGUE, ET C'EST VOULU. Le pendjabi
%  s'ecrit en shahmoukhi au Pakistan et en gourmoukhi en Inde ;
%  Ethnologue ne compte que l'occidental (« pnb »). Ce ne sont pas
%  deux graphies d'un meme texte a produire par machine : les deux
%  usages ont chacun leur lexique et leurs emprunts, et la colonne
%  « pa » se traduit pour elle-meme. Ici, la serie est un
%  « سلسلہ » ; la-bas, ce sera une « ਲੜੀ ».
%
%  UNE LANGUE QUI S'ECRIT DE DROITE A GAUCHE. Le fichier, lui, se lit
%  dans l'ordre logique — le premier caractere d'un bloc est le
%  premier mot de la phrase — et c'est le navigateur qui retourne
%  l'affichage : html.py pose « dir="rtl" » sur la seule case
%  pendjabie. Les renvois restent donc ecrits
%  \textsuperscript{(18)}, avec leurs chiffres latins, APRES le
%  substantif comme dans les neuf autres colonnes ; l'algorithme
%  bidirectionnel les remet du bon cote tout seul. Ne rien inverser a
%  la main : ce qui est retourne deux fois revient a l'endroit.
%
%  LES CHIFFRES RESTENT LATINS, renvois et numeros d'alinea. Le
%  Pakistan imprime volontiers ۱۲۳ ; mais ces chiffres-la numerotent
%  la gravure, et doivent se lire d'une colonne a l'autre. C'est le
%  parti du hindi, du bengali et des sept autres.
%
%  LE PENDJABI POSTPOSE SES RAPPORTS et place le verbe a la fin,
%  comme le hindi : on garde l'ordre des RENVOIS de l'ido, et l'on
%  refait la phrase autour.
%
%  LE NOM DE L'EDITEUR SE PRONONCE « del-MASS ». Delmas est un nom du
%  Midi -- « del mas », celui du mas -- et son s final se fait
%  entendre : on ecrit دلماس, non دلما.
% ===================================================================

%%K t01-apar-0 apar
\VUsaut{2.0mm}
\VUcentre{12.6pt}{120}{وضاحتی کتابچہ}
\VUsaut{3.2mm}
\VUcentre{8.4pt}{160}{نال لگیا}
\VUsaut{2.6mm}
\VUtitre{15.6pt}{0}{دلماس دے مددگار نقشیاں نال}
\VUsaut{3.4mm}
\VUornamento{9pt}
\VUsaut{5.0mm}
\VUcentre{13.2pt}{90}{پہلا سلسلہ}
\VUsaut{3.0mm}
\VUfilet{16mm}
\VUsaut{5.6mm}

%%K t01-apar-1 apar
\VUtitre{15.0pt}{60}{نقشہ نمبر 1}
\VUsaut{1.4mm}
\VUtitre{11.4pt}{0}{سکول۔ --- ثانوی سکول۔ --- جماعت۔}
\VUsaut{2.2mm}
\VUfilet{20mm}
\VUsaut{6.4mm}

%%K t01-tit-1 sub
\VUpk{10.2pt}{40}{عام وضاحت۔}
\VUsaut{3.0mm}

%%K t01-01-1 p
1. --- ایس نقشے وچ اک \VUgras{ثانوی سکول} دی \VUgras{جماعت} وکھائی گئی اے۔ ایس
جماعت وچ اٹھ \VUgras{میزاں} \textsuperscript{(18)} تے اٹھ \VUgras{بینچ}
\textsuperscript{(32)} نیں۔ ہر بینچ اُتے تن طالب علم بہندے نیں۔ \VUgras{ماسٹر}
\textsuperscript{(7)} \VUgras{کرسی} \textsuperscript{(8)} اُتے اپݨے
\VUgras{چبوترے} \textsuperscript{(6)} دے پچھے بیٹھے نیں۔

%%K t01-02-1 p
2. --- چبوترے دے کھبے پاسے شیشے دی اک \VUgras{الماری} \textsuperscript{(15)} کھڑی
اے۔ سجے پاسے \VUgras{کالا تختہ} \textsuperscript{(1)} اے، تے اوس اُتے
\VUgras{جھاڑن} \textsuperscript{(2)} تے \VUgras{چاک} \textsuperscript{(3)}۔

%%K t01-02-2 p
چبوترے دے اُتے \VUgras{دیواری تصویراں} \textsuperscript{(9, 11, 12)} تے اک
\VUgras{نقشہ} \textsuperscript{(10)} ٹنگے ہوئے نیں۔

%%K t01-03-1 p
3. --- سارے طالب علم اپݨی \VUgras{جگہ} \textsuperscript{(17)} اُتے نئیں۔ جان
\textsuperscript{(4)} کالے تختے کول کھڑا اے؛ اوہ جھاڑن پھڑی ہوئی اے تے
\VUgras{ہندسی شکلاں} مٹا رہیا اے۔

%%K t01-03-2 p
پیٹر چبوترے کول اے؛ اوہ \VUgras{لکھتی مشقاں} \textsuperscript{(5)} کٹھیاں کر کے
ماسٹر کول لے جا رہیا اے۔

%%K t01-04-1 p
4. --- ایہہ ماسٹر \VUgras{نویں زباناں} دے ماسٹر نیں۔ اوہ طالب علماں نوں پرائیاں
زباناں \VUgras{(جرمن، انگریزی، ہسپانوی، اطالوی، روسی} یا \VUgras{فرانسیسی)}
بولݨا، پڑھݨا تے لکھݨا سکھاندے نیں۔

%%K t01-05-1 p
5. --- جماعت بہت وڈی تے بہت روشن اے۔ سرے اُتے ہوا تے چاݨن آوݨ لئی دو
\VUgras{روشندانواں} \textsuperscript{(78)} والی اک چوڑی \VUgras{کھڑکی} تے اک
\VUgras{شیشے دا دروازہ} اے۔

%%K t01-05-2 p
ایہہ جماعت ٹھنڈی نئیں۔ وچکار اوہنوں نیم رکھݨ لئی اک بہت چنگی
\VUgras{انگیٹھی} \textsuperscript{(46)} رکھی ہوئی اے، اپݨی \VUgras{نالی}
\textsuperscript{(45)} سمیت۔

%%K t01-05-3 p
کندھ اُتے \VUgras{کھونٹیاں} \textsuperscript{(58)} گڈیاں ہوئیاں نیں؛ اوہناں اُتے
طالب علماں دیاں ٹوپیاں تے کوٹ ٹنگے رہندے نیں۔

%%K t01-tit-2 sub
\VUsaut{5.2mm}
\VUpk{10.2pt}{40}{کھیڈ دا وِہڑا۔}
\VUsaut{3.6mm}

%%K t01-06-1 p
6. --- \VUgras{گھڑی} \textsuperscript{(63)} نو وج کے پنج منٹ وکھا رہی اے۔

%%K t01-06-2 p
\VUgras{چھٹی} \textsuperscript{(76)} ہُݨے مکی اے تے سبق مُڑ شروع ہو رہیا اے۔ چھٹی
\VUgras{وِہڑے} \textsuperscript{(75)} وچ ہوندی اے۔ سانوں کجھ طالب علم کھیڈدے
وکھائی دیندے نیں۔ اک \VUgras{نگران ماسٹر} \textsuperscript{(74)} اوہناں اُتے نظر
رکھی ہوئی اے۔

%%K t01-07-1 p
7. --- وِہڑے وچ سانوں ہور وی کئی لوک وکھائی دیندے نیں : \VUgras{معائنہ کار}
\textsuperscript{(73)}، \VUgras{مہتمم} \textsuperscript{(71)} تے \VUgras{ناظم}
\textsuperscript{(72)}۔

%%K t01-07-2 p
معائنہ کار سکولاں دا معائنہ کردے نیں، مہتمم سکول چلاندے نیں، ناظم طالب علماں دی
پڑھائی اُتے نظر رکھدے نیں۔

%%K t01-07-3 p
چوتھے بندے \VUgras{چپڑاسی} \textsuperscript{(77)} نیں۔ اوہ غیر حاضر طالب علماں دے
ناں لکھݨ لئی بغل وچ اک رجسٹر دبائی ہوئے نیں۔

%%K t01-08-1 p
8. --- وِہڑے دے سرے اُتے \VUgras{ورزش دے سامان} \textsuperscript{(70)} تے
\VUgras{دستکاری} \textsuperscript{(69)} دا کارخانہ اے۔

%%K t01-08-2 p
نقشے دے سجے پاسے \VUgras{موسیقی} تے \VUgras{مصوری} دے \VUgras{کمرے} مشکل نال
وکھائی دیندے نیں۔

%%K t01-noto-1 noto
\VUnotes{143.2mm}{%
(*) \textit{ناواں} لئی وی ایہی صلاح دِتی جاندی اے پئی اوہناں نوں اوہناں دی
لاطینی شکل وچ لکھیا جاوے۔ بے شمار شکلاں توں بچݨ دا ایہو اک ای طریقہ اے۔ فیر،
\textit{Giovanni, Jean, Johann, Jan, Hans, John, Ivan} وچوں چوݨ وی کیویں؟ کیہ
ناواں دا وکھرا کوش بݨایا جاوے؟ لاطینی فاعلی شکل لے لؤ تے آکھو :
\VUgras{Ioannes, Ioanna; Iulius, Iulia; Iakobus; Andreas; Lukas.}
(مکمل قواعد)۔%
}

%%K t01-tit-3 sub
\VUpk{10.2pt}{40}{تفصیلی وضاحت۔}
\VUsaut{2.4mm}

%%K t01-tit-4 sub
\VUpk{10.2pt}{40}{چنگے طالب علم۔}
\VUsaut{3.0mm}

%%K t01-09-1 p
9. --- چبوترے دے سجے پاسے پہلی بینچ اُتے بیٹھے طالب علم \VUgras{ہنری} (*)،
\VUgras{سٹیفن} تے \VUgras{چارلس} نیں۔

%%K t01-09-2 p
ہنری بہت دھیان دیݨ والا تے محنتی اے؛ اوہ ماسٹر دی گل سُݨ رہیا اے۔ اوہدی میز
اُتے اک \VUgras{کاپی} \textsuperscript{(28)}، اک \VUgras{کتاب}
\textsuperscript{(29)}، اک \VUgras{لغت} \textsuperscript{(30)} تے اک
\VUgras{قلمدان} \textsuperscript{(31)} اے۔ اوہدی کتاب وچ بہت سارے \VUgras{ورقے}
نیں؛ ہر باب وچ کئی \VUgras{پیرے} نیں تے ہر \VUgras{سطر} وچ بہت سارے
\VUgras{لفظ}۔

%%K t01-09-4 p
اوہدا گوانڈھی سٹیفن \textsuperscript{(24)} جوشیلا اے۔ اوہ اگلی واری دا دِتا ہویا
سبق اپݨی کاپی وچ لکھ رہیا اے۔ اوہدی \VUgras{لکھائی} سوہݨی اے۔ اوہدی کتاب بند
اے۔ نِرا \VUgras{جِلد} \textsuperscript{(27)} وکھائی دیندی اے۔ اوہدے کول اوہدا
\VUgras{پرکار} \textsuperscript{(26)} پیا اے۔

%%K t01-09-5 p
چارلس \textsuperscript{(23)} ہتھ وچ کتاب پھڑی ہوئی اے؛ اوہ اوہنوں کھول کے اپݨا
سبق مُڑ دہرا رہیا اے۔ ہر طالب علم وانگوں اوہدے سامݨے وی اک \VUgras{دوات}
\textsuperscript{(25)} اے۔ اوہ فرمانبردار اے۔

%%K t01-10-1 p
10. --- اگلی میز اُتے \VUgras{لوئی، انٹونی} تے \VUgras{پال} نیں۔ لوئی اپݨا
\VUgras{سبق} \textsuperscript{(22)} سُݨا رہیا اے تے ماسٹر دے \VUgras{سوالاں} دا
جواب دے رہیا اے۔ انٹونی \textsuperscript{(21)} بہت بے دھیان اے؛ ماسٹر کدے کدے
اوہنوں سزا وی دیندے نیں۔ پال \textsuperscript{(20)} چنگا یار اے، ملݨسار تے
مددگار۔ اوہ بہت دھیان دیندا اے۔

%%K t01-tit-5 sub
\VUsaut{4.2mm}
\VUpk{10.2pt}{40}{مندے طالب علم۔}
\VUsaut{3.2mm}

%%K t01-11-1 p
11. --- ہنری دے پچھے \VUgras{جارج} \textsuperscript{(38)} بیٹھا اے۔ اوہ مندا مُنڈا
نئیں، پر \VUgras{باتونی}، بے دھیان تے شرارتی اے۔ اوہدی \VUgras{چنچلتا} تے
\VUgras{نافرمانی} سبق دے ویلے \VUgras{خلل} پاندیاں نیں، تے اوہنوں اکثر سخت
\VUgras{تنبیہ} ملݨی چاہیدی اے۔ اوہدی \VUgras{کچی کاپی} \textsuperscript{(35)}
گندی اے، اوہ اوہدے اُتے اپݨے \VUgras{قلم} \textsuperscript{(34)} نال
\VUgras{سیاہی دے دھبے} پاندا رہندا اے، ایسے لئی اوہنوں ہمیشہ اپݨا \VUgras{ربڑ}
\textsuperscript{(36)} چاہیدا اے؛ اوہدا \VUgras{بستہ} \textsuperscript{(33)} پھٹیا
ہویا اے، کیوں جے اوہ بہت لاپرواہ اے۔ نویں زباناں دی جماعت وچ اوہ اپݨا
\VUgras{پنسل دان} \textsuperscript{(37)} کیوں لیاندا اے؟

%%K t01-11-2 p
اوہدا گوانڈھی ایڈمنڈ \textsuperscript{(39)} اوہنوں پسند نئیں کردا۔ اوہ منّیے تے
تھوڑا سست اے، پر چپ چپیتا تے ٹھہریا ہویا اے۔ اوہ بک بک نئیں کردا؛ بس، سُݨن دی
تھاں اوہ اپݨی \VUgras{درسی کتاب} دیاں \VUgras{کہاݨیاں} پڑھݨا ودھ پسند کردا اے۔

%%K t01-11-3 p
ایس بینچ دا تیجا طالب علم \VUgras{لڈووک} \textsuperscript{(40)} اے۔ اوہ محنتی اے۔
اوہ \VUgras{کاغذ دے} اک چٹے \VUgras{ورقے} اُتے لکھ رہیا اے۔ اپݨے سامݨے اوہنے
اپݨا \VUgras{سوکھتہ} \textsuperscript{(41)} رکھیا ہویا اے۔

%%K t01-12-1 p
12. --- ماسٹر دے کھبے پاسے دوجی بینچ دے طالب علم بہت نِکے نیں۔ پہلا، نِکا
\VUgras{ایمل}، ہالے چنگی طرحاں لکھ نئیں سکدا؛ اوہ اپݨی \VUgras{سلیٹ}
\textsuperscript{(42)}، اپݨی \VUgras{سلیٹی} تے اپݨا \VUgras{اسفنج}
\textsuperscript{(43)} لیاندا اے۔

%%K t01-12-2 p
اوہدے کول \VUgras{اگست} تے \VUgras{برنارڈ} \textsuperscript{(44)} نیں۔ پچھلا چنگا
کم کردا اے، پر اوہدا \VUgras{لباس} بہت میلا کچیلا رہندا اے۔

%%K t01-13-1 p
13. --- اوہناں دے پچھے تن \VUgras{سست} بیٹھے نیں۔ \VUgras{البرٹ} اپݨے سبق یاد
نئیں کردا۔ \VUgras{جیمز} \textsuperscript{(47)} سُݨن دی تھاں سُتا پیا اے۔ اوہدی
\VUgras{سستی} اوہنوں \VUgras{جھاڑ} تے \VUgras{سزا} تیکر وی پہنچاندی اے۔ چھیکر
\VUgras{کرسچن} \textsuperscript{(48)} کِتے ودھ شرارتی اے۔ اوہدا \VUgras{چال چلن}
مندا اے تے اوہ سُدھرݨ دا کوئی جتن نئیں کردا۔

%%K t01-14-1 p
14. --- اوہناں دے گوانڈھی، ایہدے اُلٹ، سُشیل نیں۔ \VUgras{ارنسٹ}
\textsuperscript{(51)} نوں ترتیب تے قاعدہ پیارے نیں؛ اوہ ہمیشہ اپݨا
\VUgras{پیمانہ} \textsuperscript{(50)} تے اپݨی \VUgras{پنسل} \textsuperscript{(49)}
کم وچ لیاندا اے۔ اوہ \VUgras{باہرلا طالب علم} اے۔

%%K t01-14-2 p
\VUgras{فرانسس} \textsuperscript{(52)} \VUgras{ہاسٹل والا} اے؛ اوہ وردی پاندا اے،
تے سکول وچ ای کھاندا تے سوندا اے۔ اوہ چنگا \VUgras{یار} اے۔

%%K t01-14-3 p
مورس اپݨے \VUgras{چاقو} \textsuperscript{(56)} نال اپݨی \VUgras{پنسل} چھِل رہیا
اے؛ اوہ بہت ہوشیار اے۔

%%K t01-14-4 p
اوہ اپݨیاں ساریاں کتاباں لیاندا اے، اپݨے \VUgras{قواعد} \textsuperscript{(55)}،
اپݨی \VUgras{نوٹ بک} \textsuperscript{(54)}، تے اک \VUgras{لکھݨ دا گدا}
\textsuperscript{(53)} وی۔ اوہدا \VUgras{تھیلا} \textsuperscript{(57)} اوہدے پچھے
تھلے پیا اے۔

%%K t01-14-5 p
جماعت دے پچھلے پاسے دیاں دو میزاں اُتے \VUgras{چنگے طالب علم}
\textsuperscript{(19)} بیٹھے نیں۔

%%K t01-tit-6 sub
\VUsaut{4.6mm}
\VUpk{10.2pt}{40}{مصوری تے موسیقی۔}
\VUsaut{3.4mm}

%%K t01-15-1 p
15. --- \VUgras{مصوری دے کمرے} وچ کندھ اُتے سوہݨے \VUgras{نمونے} ٹنگے نیں تے
چھت توں اک \VUgras{چھجے} سمیت \VUgras{لیمپ} \textsuperscript{(62)} لٹکیا ہویا اے۔
\VUgras{ماسٹر} \textsuperscript{(60)} بہت صابر نیں؛ اوہ طالب علماں دیاں
\VUgras{تصویراں} \textsuperscript{(61)} سُدھاردے نیں تے اوہناں نوں سامݨے رکھی
\VUgras{مورتی} \textsuperscript{(59)} ویکھ کے تصویر بݨاؤݨا سکھاندے نیں۔

%%K t01-16-1 p
16. --- \VUgras{موسیقی دا کمرہ} بہت دلچسپ لگدا اے۔ اوتھے \VUgras{موسیقی} پڑھݨا،
\VUgras{سُر لپی} اُتے لکھے \VUgras{سُر} \textsuperscript{(67)} گاؤݨا (ڈو، رے، می،
فا، سول، لا، سی\,=\,C. D. E. F. G. A. B.)، \VUgras{کنجیاں} پچھاݨنا تے سوہݨیاں
\VUgras{دھناں} گاؤݨا سکھایا جاندا اے۔ سُر لپیاں \VUgras{رحل} \textsuperscript{(65)}
اُتے رکھیاں جاندیاں نیں۔ اک طالب علم \VUgras{پیانو وجا رہیا} \textsuperscript{(64)}
اے تے \VUgras{موسیقی دے ماسٹر} \textsuperscript{(66)} \VUgras{تال دے رہے}
\textsuperscript{(68)} نیں۔

%%K t01-17-1 p
17. --- سانوں \VUgras{دستکاری} \textsuperscript{(69)} دے \VUgras{کارخانے} دا اک
کونا مشکل نال وکھائی دیندا اے، جِتھے جواک لکڑ چھِلݨا تے لوہا ریتݨا سکھ سکدے
نیں۔ ایہہ ہر سکول وچ نئیں ہوندا۔

%%K t01-tit-7 sub
\VUsaut{5.0mm}
\VUpk{10.2pt}{40}{سائنس دا کمرہ۔}
\VUsaut{3.6mm}

%%K t01-18-1 p
18. --- ریاضی دے ماسٹر طالب علماں نوں \VUgras{ہندسہ، حساب، گݨتی} تے
\VUgras{میٹری نظام} پڑھاندے نیں۔ نقشے اُتے سانوں ہندسے دیاں مکھ شکلاں وکھائی
دیندیاں نیں : اک \VUgras{دائرہ} \textsuperscript{(a)}، اک \VUgras{مربع}
\textsuperscript{(b)}، اک \VUgras{مثلث} \textsuperscript{(c)}، اک \VUgras{لکیر}
\textsuperscript{(d)}۔

%%K t01-18-2 p
سانوں اک \VUgras{حساب} وی وکھائی دیندا اے؛ اوہ نہ \VUgras{جمع} اے، نہ
\VUgras{تفریق}، نہ \VUgras{ضرب} : اوہ \VUgras{تقسیم} \textsuperscript{(e)} اے۔

%%K t01-19-1 p
19. --- میٹری نظام دے نقشے اُتے سانوں وکھائی دیندے نیں : اک \VUgras{میٹر}
\textsuperscript{(a)}، اک \VUgras{ترازو} \textsuperscript{(b)} تے اوہدے
\VUgras{وٹے}، اک \VUgras{لیٹر} \textsuperscript{(e)}، اک \VUgras{ڈیکالیٹر}
\textsuperscript{(f)}، اک \VUgras{ڈیسی لیٹر} تے اک \VUgras{مکعب میٹر}
\textsuperscript{(d)}۔

%%K t01-19-2 p
طالب علم \VUgras{قدرتی علوم} \textsuperscript{(11)} --- حیوانیات
\textsuperscript{(a)}، نباتیات \textsuperscript{(b)}، ارضیات \textsuperscript{(c)}
--- تے \VUgras{تاریخ} \textsuperscript{(12)} وی پڑھدے نیں۔

%%K t01-19-3 p
الماری وچ \VUgras{طبیعیات} \textsuperscript{(15)} تے \VUgras{کیمیا}
\textsuperscript{(16)} دے سامان رکھے نیں۔

%%K t01-19-4 p
الماری دے اُتے نیں : اک \VUgras{گولا} \textsuperscript{(14)}، اک \VUgras{مخروط}
تے اک \VUgras{کرہ ارض} \textsuperscript{(13)}۔

%%K t01-19-5 p
نقشیاں دے اُتے اک \VUgras{نقشہ} ٹنگیا ہویا اے جیدے اُتے دنیا دے پنجے حصے وکھائے
گئے نیں : \VUgras{یورپ} \textsuperscript{(g)}، \VUgras{ایشیا}
\textsuperscript{(e)}، \VUgras{افریقہ} \textsuperscript{(f)}، \VUgras{امریکہ}
\textsuperscript{(ab)} تے \VUgras{اوشیانیا} \textsuperscript{(h)}، تے دو وڈے
سمندر، \VUgras{بحرالکاہل} \textsuperscript{(c)} تے \VUgras{بحر اوقیانوس}
\textsuperscript{(d)}۔
\VUsaut{6.0mm}
\VUfilet{22mm}
